\section{Raw comparison}\label{sec:raw}

Do paragliders and hang gliders have the same $H$ when all terrain is taken together?

\paragraph{A, fixed population.} With $T=$ \UniTauFixed{} s the population is the \UniNFirstAllAfixed{} paraglider flights
(\UniPFirstAllAfixed{} pilots) and the \UniNSecondAllAfixed{} hang glider flights (\UniPSecondAllAfixed{} pilots) that have a
piece of \UniPieceFixed{} s. The exponent is $H=$ \UniHFirstAllAfixed{} \UniHBandFirstAllAfixed{} for paragliders and
\UniHSecondAllAfixed{} \UniHBandSecondAllAfixed{} for hang gliders; the difference is $\Hpara-\Hhang=$ \UniDHAllAfixed{}
with the band \UniDHBandAllAfixed{} (Table~\ref{tab:raw}). Fig.~\ref{fig:raw-A}(a) shows why: the curve of the paragliders lies
above the other at every lag, by \UniDAlphaAllAfixedHundred{} in $\alpha$ at 100 s and \UniDAlphaAllAfixedThousand{} at 1000 s.
Both curves rise to a maximum at some tens of seconds and fall to a plateau.

With the longest $T$, \UniTAllAlongest{} s, the populations are \UniNFirstAllAlongest{} paraglider and \UniNSecondAllAlongest{} hang
glider flights (Fig.~\ref{fig:raw-A}(b)): $H=$ \UniHFirstAllAlongest{} and \UniHSecondAllAlongest, a difference of
\UniDHAllAlongest{} with the band \UniDHBandAllAlongest, which includes 0. Above 1000 s the curves cross and the band of the
difference of $\alpha$ is wide, since only \UniNSecondAllAlongest{} hang glider flights are left.

\begin{figure}
\centering
\includegraphics[width=\textwidth]{raw_A.pdf}
\caption{On a fixed population, $\alpha(\tau)$ of paragliders lies above that of hang gliders up to 1000 s, and the difference
with the longest $T$ is no longer distinguishable from 0 above that lag. All terrain, design A, (a) $T=$ \UniTauFixed{} s, (b) the
longest $T$. Top: $\alpha(\tau)$ with the dashed $2H$ of each fit; bottom: $\Delta\alpha=\alpha_{\rm para}-\alpha_{\rm hang}$ with the dashed $2(\Hpara-\Hhang)$.
Bands 90\%; numbers of flights above the panels.}\label{fig:raw-A}
\end{figure}

\paragraph{B, all flights.} Taking all the flights each lag can use, $H$ is \UniHFirstAllBkilo{} and \UniHSecondAllBkilo{}
over the lags up to $1000$ s, and \UniHFirstAllBlongest{} and \UniHSecondAllBlongest{} over the lags up to \UniTAllBlongest{} s;
the differences are \UniDHAllBkilo{} \UniDHBandAllBkilo{} and \UniDHAllBlongest{} \UniDHBandAllBlongest{}. Fig.~\ref{fig:raw-B} shows
the population: at long lags the number of flights falls and $\alpha(\tau)$ of both disciplines bends upward, which a fixed
population does not do. Section~\ref{sec:fixed} compares the two designs.

\begin{figure}
\centering
\includegraphics[width=\textwidth]{raw_B.pdf}
\caption{When each lag takes all the flights it can, the flights fall at long lags and $\alpha(\tau)$ bends upward there. All terrain,
design B; same panels as Fig.~\ref{fig:raw-A} and, bottom, the number of flights at each lag.}\label{fig:raw-B}
\end{figure}

\begin{table}
\centering
\small
\renewcommand{\arraystretch}{1.15}\setlength{\tabcolsep}{4pt}
\begin{adjustbox}{max width=\textwidth}
\begin{tabular}{lrrccc}
\toprule
Population & flights (pilots) para & flights (pilots) hang & $H$ para & $H$ hang & $\Delta H$ \\
\midrule
all terrain, A, $T=1000$ s & 55\,074 (5\,505) & 1\,400 (189) & \shortstack{0.922\\[-1pt]{\scriptsize [0.921, 0.923]}} & \shortstack{0.898\\[-1pt]{\scriptsize [0.890, 0.905]}} & \shortstack{$+$0.024\\[-1pt]{\scriptsize [$+$0.017, $+$0.031]}} \\
all terrain, A, longest $T$ (11450 s) & 7\,599 (2\,074) & 172 (60) & \shortstack{0.923\\[-1pt]{\scriptsize [0.920, 0.926]}} & \shortstack{0.915\\[-1pt]{\scriptsize [0.902, 0.928]}} & \shortstack{$+$0.007\\[-1pt]{\scriptsize [$-$0.006, $+$0.020]}} \\
all terrain, B, to 1000 s & 55\,074 (5\,505) & 1\,400 (189) & \shortstack{0.922\\[-1pt]{\scriptsize [0.921, 0.923]}} & \shortstack{0.897\\[-1pt]{\scriptsize [0.890, 0.905]}} & \shortstack{$+$0.024\\[-1pt]{\scriptsize [$+$0.017, $+$0.032]}} \\
all terrain, B, to longest $T$ (11450 s) & 7\,599 (2\,074) & 172 (60) & \shortstack{0.917\\[-1pt]{\scriptsize [0.914, 0.920]}} & \shortstack{0.897\\[-1pt]{\scriptsize [0.883, 0.911]}} & \shortstack{$+$0.020\\[-1pt]{\scriptsize [$+$0.006, $+$0.033]}} \\
\bottomrule
\end{tabular}
\end{adjustbox}
\caption{H of paragliders and hang gliders, all terrain: design A (fixed population, $T$ shown) and B (all flights); flights (pilots) at the longest lag of the design; bands 90\%, difference para $-$ hang.}\label{tab:raw}
\end{table}


\keep{All terrain, on the fixed population: $\Hpara-\Hhang=$ \UniDHAllAfixed{} \UniDHBandAllAfixed{} at $T=$ \UniTauFixed{} s,
\UniDHAllAlongest{} \UniDHBandAllAlongest{} at $T=$ \UniTAllAlongest{} s. Paragliders have the larger $H$ at 1000 s; the difference
is small, and at the longest $T$ it is not resolved.}
